\documentclass{article}
\usepackage{import}
\import{../../template/}{article-template}
\graphicspath{{../../images/}}
\addbibresource{../../template/library.bib}

\title{Fluent}
\author{PENG}
\date{Created: 20260917, Version: 20260917}

\begin{document}
\maketitle
\tableofcontents
\section{Introduction}
\href{https://ansyshelp.ansys.com/public/account/secured?returnurl=/Views/Secured/corp/v261/en/flu_ug/flu_ug.html}{Fluent User's Guide}

\begin{definition}[Ansys Fluent]
    Ansys Fluent a computer program written in C for modelling fluid flow, heat transfer, and chemical reactions in complex geometries.
\end{definition}

\section{Basic Steps for CFD Analysis}
\begin{enumerate}
    \item define the modeling goals:
          simplifying assumptions, boundary conditions
    \item create the model geometry and mesh:
          \begin{itemize}
              \item unstructured meshes, default units: meters
              \item mesh resolution sufficient in each region of the domain
              \item mesh check: cell skewness $< 0.98$
          \end{itemize}
    \item set up the solver and physical models
          \begin{itemize}
              \item boundary conditions at all boundary zones
              \item initial solution
              \item converged solution is not necessarily a correct one
              \item use \textbf{second-order upwind discretization} for better accuracy
          \end{itemize}
    \item compute and monitor the solution
          \begin{itemize}
              \item number of iterations to reach a converged solution
          \end{itemize}
    \item examine and save the results
          \begin{itemize}
              \item flow features: flow pattern, separation, .etc
              \item quantitative results: forces and moments
          \end{itemize}
    \item revisions to the numerical or physical model parameters
\end{enumerate}

\section{Executing Ansys Fluent}

\end{document}